% =============================================================================
% DRAFT — NOT INCLUDED IN THE BUILD
% =============================================================================
% Purpose: proposed expansion of Round 2, Step 4 of the security proof
% (sections/security.tex), as requested: discuss the reprogramming of the
% random oracle when simulating the NIZK proof pi_simid, and the associated
% bad event (reprogramming failure / collision with a prior RO query).
%
% WHAT WE WANTED TO CHANGE
% ------------------------
% In sections/security.tex, Round 2 currently has a single terse item for
% Step 4:
%
%   \item Generate a simulated an online extractable zero-knowledge proof
%         $\pi_\simid$ (following Figure~\ref{fg:concrete-simulator}) for the
%         relation $\mathcal{R}_P$ w.r.t.\@ $C_\simid[0]$.
%
% The plan was to (a) replace that line with the detailed Fiat--Shamir / Schnorr
% HVZK simulation below (which makes the RO reprogramming explicit), and
% (b) add the "Bad event: reprogramming failure" paragraph after the
% \end{itemize} of Round 2, before the "Round 3" paragraph.
%
% CAVEATS TO REVIEW BEFORE INTEGRATING
% ------------------------------------
%  - The R_h formula assumes h^{\sk} is a public group element available to the
%    verifier in equation (3) of R_P. The relation R_P in prelim.tex still has
%    \simon{} TODO notes ("Verify the need of (1) and define what is c'"), so the
%    exact statement/verifier inputs may change. Adjust the commitment formulas
%    to whatever the finalized R_P verification equations require.
%  - $\badevent$ and $q_h$ are used as ad-hoc notation here; define/rename to
%    match the rest of the paper if needed.
%  - The bound assumes a single simulated proof per execution (only party
%    $\simid$ is simulated). If multiple proofs are simulated, add a union bound.
% =============================================================================


% ---- (a) Replacement for the Step-4 \item ------------------------------------

        \item Generate a simulated online-extractable zero-knowledge proof $\pi_\simid$ for the
          statement $C_\simid[0]$ with context $\gamma_\simid \gets (\sid, \simid, C_\simid[0])$,
          following the simulator of Figure~\ref{fg:concrete-simulator} for the relation $\mathcal{R}_P$.
          Because $C_\simid[0]$ is defined in Equation~\eqref{eq:simid-comm} so as to embed the DKG
          challenge $\dkgchal$, $\SIM$ does \emph{not} know the corresponding witness, i.e.\@ the
          discrete logarithm $\dl$ of $C_\simid[0]$. It therefore cannot run the honest prover and
          instead invokes the honest-verifier zero-knowledge simulator of the underlying Schnorr
          proof, reprogramming the random oracle $\HashSchnorrR$ to make the transcript verify:
          \begin{enumerate}
            \item Form the ElGamal ciphertext $c \gets (C_0, C_1)$ honestly, sampling fresh randomness
              $r \rgets \Zq$ and setting $C_0 \gets g^{r}$ and $C_1 \gets h^{r} \cdot C_\simid[0]$.
              Note that this uses only the public value $g^{\sk} = C_\simid[0]$ and not $\sk$ itself.
            \item Sample a challenge $c' \rgets \Zq$ and a response $z \rgets \Zq$ uniformly and
              independently at random.
            \item Compute the commitments $R_g$ and $R_h$ by rearranging the verification equations~(2)
              and~(3) of $\mathcal{R}_P$, i.e.\@ $R_g \gets g^{z} \cdot (g^{\sk})^{-c'}$ and
              $R_h \gets h^{z} \cdot (h^{\sk})^{-c'}$, where $g^{\sk} = C_\simid[0]$ and $h^{\sk}$ are
              the public group elements that the verifier uses in those equations.
            \item Reprogram the random oracle so that $\HashSchnorrR(\gamma_\simid, c, R_g, R_h) \gets c'$,
              and set $\pi_\simid \gets (c, z, R_g, R_h)$.
          \end{enumerate}


% ---- (b) New paragraph to insert after Round 2's \end{itemize} ---------------

    \medskip

    \paragraph{\bf Bad event: reprogramming failure.}
    The last step of the simulation above overwrites the value of $\HashSchnorrR$ at the point
    $(\gamma_\simid, c, R_g, R_h)$. This reprogramming is sound only if the adversary has \emph{not}
    already queried $\HashSchnorrR$ at exactly this point: otherwise $\SIM$ would be forced to assign
    two distinct outputs to the same input, and the resulting inconsistency could be detected by
    $\adv$. We denote by $\badevent$ the event that this collision occurs, and we let $\SIM$ abort
    whenever $\badevent$ is triggered.

    We bound $\Pr[\badevent]$. The commitment $R_g = g^{z} \cdot C_\simid[0]^{-c'}$ is a uniformly
    random element of $\Gr$, since the response $z$ is sampled uniformly and independently of the
    adversary's view prior to the production of $\pi_\simid$; consequently the reprogramming point
    $(\gamma_\simid, c, R_g, R_h)$ is unpredictable to $\adv$. Letting $q_h$ denote the number of
    queries that $\adv$ makes to $\HashSchnorrR$, the probability that the (single) simulated proof
    collides with one of these queries is therefore at most $q_h / \lvert \Gr \rvert = q_h / q$,
    which is negligible in $\lambda$ since $q$ is exponential in $\lambda$.

    Conditioned on $\badevent$ not occurring, the reprogramming is undetectable, and by the perfect
    honest-verifier zero-knowledge property of the Schnorr proof system the simulated proof
    $\pi_\simid$ is distributed identically to one produced by an honest prover holding the witness
    $\dl$. Hence the adversary's view in this hybrid is statistically indistinguishable from its view
    in the real protocol, the statistical distance being bounded by $q_h / q$.
